```latex
\subsection{From the frequency law to novel continuations}
\label{sec:novel-continuations}

We call a validation target \emph{novel} if its context occurs in training but was never followed by that target there. The word itself may be familiar; novelty is a property of the context--target pair. Let $f(F)$ denote the fraction of validation positions in context-count bin $F \ge 1$ whose targets are novel.

Splitting the gap by novelty gives the approximation
\begin{equation}
\mathrm{gap}_t(F) \;\approx\; a_t\, f(F) + b_t\,[1 - f(F)],
\label{eq:coverage-main}
\end{equation}
where $a_t$ and $b_t$ are the typical loss differences for novel and seen targets. The exact $F$-dependent decomposition and correction terms appear in Appendix~\ref{app:math}.

Empirically, $f(F)$ declines roughly as $F^{-\alpha}$, a corpus statistic that conditional heavy tails and Good--Turing estimation help explain. When the novel term dominates and the component differences vary slowly, the gap inherits approximately the same decay, though the exponents need not match exactly. Figure~\ref{fig:frequency} overlays the three-seed mean online gap with an affine fit to the measured $f(F)$; the fitted coefficients are effective values rather than literal component means (fitting details in Appendix~\ref{app:math}). Coverage thus explains the frequency law, but not why replay harms these novel continuations---the question we turn to next.
```